\documentclass[11pt]{article}
\newcommand{\ListaTitulo}{Lista 02: Tendência central, quantis e boxplot}
% Preâmbulo compartilhado: MATF14, Listas de Exercícios 2026
% Usado por ListaNN.tex e GabaritoNN.tex via % Preâmbulo compartilhado: MATF14, Listas de Exercícios 2026
% Usado por ListaNN.tex e GabaritoNN.tex via % Preâmbulo compartilhado: MATF14, Listas de Exercícios 2026
% Usado por ListaNN.tex e GabaritoNN.tex via \input{preamble}

\usepackage[utf8]{inputenc}
\usepackage[T1]{fontenc}
\usepackage[brazil]{babel}
\usepackage{fullpage}
\usepackage{amsmath,amssymb}
\usepackage{enumitem}
\usepackage{xcolor}
\usepackage{fancyhdr}
\usepackage{titlesec}
\usepackage{listings}
\usepackage{hyperref}
\usepackage{booktabs}
\usepackage{graphicx}
\usepackage{tikz}

\definecolor{matf14azul}{HTML}{0D3B54}
\definecolor{matf14dourado}{HTML}{C9971E}

\titleformat{\section}{\normalfont\Large\bfseries\color{matf14azul}}{\thesection}{1em}{}
\titleformat{\subsection}{\normalfont\large\bfseries\color{matf14azul}}{\thesubsection}{1em}{}

\providecommand{\ListaTitulo}{Lista de Exercícios}

\setlength{\headheight}{14pt}
\pagestyle{fancy}
\fancyhf{}
\lhead{\small MATF14: Estatística Econômica I}
\rhead{\small \ListaTitulo}
\cfoot{\thepage}
\renewcommand{\headrulewidth}{0.4pt}

\lstset{
  basicstyle=\ttfamily\small,
  breaklines=true,
  frame=single,
  columns=fullflexible,
  backgroundcolor=\color{gray!8},
  commentstyle=\color{matf14dourado},
  keywordstyle=\color{matf14azul}\bfseries,
  language=R
}

\newcounter{questaocounter}
\newenvironment{questao}[1]{%
  \stepcounter{questaocounter}%
  \bigskip\noindent\textbf{Questão #1.}\ %
}{\par}

\newcommand{\cabecalho}[2]{%
  \begin{center}
    {\Large\bfseries\color{matf14azul} #1}\\[0.3em]
    {\large #2}
  \end{center}
  \vspace{0.5em}
  \noindent\rule{\linewidth}{0.6pt}
  \vspace{0.5em}
}

\newenvironment{solucao}{%
  \par\smallskip\noindent\textit{\color{matf14dourado!80!black}Solução.}\ %
}{\hfill$\blacksquare$\par}


\usepackage[utf8]{inputenc}
\usepackage[T1]{fontenc}
\usepackage[brazil]{babel}
\usepackage{fullpage}
\usepackage{amsmath,amssymb}
\usepackage{enumitem}
\usepackage{xcolor}
\usepackage{fancyhdr}
\usepackage{titlesec}
\usepackage{listings}
\usepackage{hyperref}
\usepackage{booktabs}
\usepackage{graphicx}
\usepackage{tikz}

\definecolor{matf14azul}{HTML}{0D3B54}
\definecolor{matf14dourado}{HTML}{C9971E}

\titleformat{\section}{\normalfont\Large\bfseries\color{matf14azul}}{\thesection}{1em}{}
\titleformat{\subsection}{\normalfont\large\bfseries\color{matf14azul}}{\thesubsection}{1em}{}

\providecommand{\ListaTitulo}{Lista de Exercícios}

\setlength{\headheight}{14pt}
\pagestyle{fancy}
\fancyhf{}
\lhead{\small MATF14: Estatística Econômica I}
\rhead{\small \ListaTitulo}
\cfoot{\thepage}
\renewcommand{\headrulewidth}{0.4pt}

\lstset{
  basicstyle=\ttfamily\small,
  breaklines=true,
  frame=single,
  columns=fullflexible,
  backgroundcolor=\color{gray!8},
  commentstyle=\color{matf14dourado},
  keywordstyle=\color{matf14azul}\bfseries,
  language=R
}

\newcounter{questaocounter}
\newenvironment{questao}[1]{%
  \stepcounter{questaocounter}%
  \bigskip\noindent\textbf{Questão #1.}\ %
}{\par}

\newcommand{\cabecalho}[2]{%
  \begin{center}
    {\Large\bfseries\color{matf14azul} #1}\\[0.3em]
    {\large #2}
  \end{center}
  \vspace{0.5em}
  \noindent\rule{\linewidth}{0.6pt}
  \vspace{0.5em}
}

\newenvironment{solucao}{%
  \par\smallskip\noindent\textit{\color{matf14dourado!80!black}Solução.}\ %
}{\hfill$\blacksquare$\par}


\usepackage[utf8]{inputenc}
\usepackage[T1]{fontenc}
\usepackage[brazil]{babel}
\usepackage{fullpage}
\usepackage{amsmath,amssymb}
\usepackage{enumitem}
\usepackage{xcolor}
\usepackage{fancyhdr}
\usepackage{titlesec}
\usepackage{listings}
\usepackage{hyperref}
\usepackage{booktabs}
\usepackage{graphicx}
\usepackage{tikz}

\definecolor{matf14azul}{HTML}{0D3B54}
\definecolor{matf14dourado}{HTML}{C9971E}

\titleformat{\section}{\normalfont\Large\bfseries\color{matf14azul}}{\thesection}{1em}{}
\titleformat{\subsection}{\normalfont\large\bfseries\color{matf14azul}}{\thesubsection}{1em}{}

\providecommand{\ListaTitulo}{Lista de Exercícios}

\setlength{\headheight}{14pt}
\pagestyle{fancy}
\fancyhf{}
\lhead{\small MATF14: Estatística Econômica I}
\rhead{\small \ListaTitulo}
\cfoot{\thepage}
\renewcommand{\headrulewidth}{0.4pt}

\lstset{
  basicstyle=\ttfamily\small,
  breaklines=true,
  frame=single,
  columns=fullflexible,
  backgroundcolor=\color{gray!8},
  commentstyle=\color{matf14dourado},
  keywordstyle=\color{matf14azul}\bfseries,
  language=R
}

\newcounter{questaocounter}
\newenvironment{questao}[1]{%
  \stepcounter{questaocounter}%
  \bigskip\noindent\textbf{Questão #1.}\ %
}{\par}

\newcommand{\cabecalho}[2]{%
  \begin{center}
    {\Large\bfseries\color{matf14azul} #1}\\[0.3em]
    {\large #2}
  \end{center}
  \vspace{0.5em}
  \noindent\rule{\linewidth}{0.6pt}
  \vspace{0.5em}
}

\newenvironment{solucao}{%
  \par\smallskip\noindent\textit{\color{matf14dourado!80!black}Solução.}\ %
}{\hfill$\blacksquare$\par}


\begin{document}

\cabecalho{Lista de Exercícios 02}{Medidas de tendência central, quantis, IQR e boxplot (Cap. 1, Aulas 04--05)}

\noindent \textit{Justifique todas as respostas. As resoluções são discutidas em sala e nos horários de atendimento; não são publicadas neste material.}

\begin{questao}{1} \textbf{(Cálculo direto: salários de uma pequena empresa)}

Os salários mensais (em R\$ mil) de 9 funcionários de uma pequena empresa são:
$$
3, \ 3, \ 3{,}5, \ 4, \ 4, \ 4{,}5, \ 5, \ 6, \ 22
$$
(o último valor é o salário do sócio-diretor).

\begin{enumerate}[label=(\alph*)]
  \item Calcule a média, a mediana e a moda dos salários.
  \item Qual medida representa melhor "o salário típico" desta empresa? Justifique com base na
    comparação entre média e mediana.
  \item Se o sócio-diretor tivesse um salário de R\$ 8 mil em vez de R\$ 22 mil, qual das três
    medidas mudaria mais? Qual mudaria menos, ou nada?
\end{enumerate}
\end{questao}

\begin{questao}{2} \textbf{(Assimetria e ordenação das medidas)}

Em uma pesquisa sobre patrimônio líquido de 1.000 famílias de um município, obteve-se
média $=$ R\$ 340 mil e mediana $=$ R\$ 210 mil.

\begin{enumerate}[label=(\alph*)]
  \item A distribuição de patrimônio é aproximadamente simétrica, assimétrica à direita ou
    assimétrica à esquerda? Justifique formalmente, usando a relação entre média e mediana vista
    em aula.
  \item Esboce (descreva em palavras a forma de) um histograma compatível com essa relação entre
    média e mediana.
  \item Um órgão de pesquisa quer comunicar "quanto tem uma família típica deste município" para o
    público em geral, em uma única frase. Qual medida você recomendaria usar, e por quê?
\end{enumerate}
\end{questao}

\begin{questao}{3} \textbf{(Quantis: tempo de espera em atendimento bancário)}

O tempo de espera (em minutos) de 12 clientes em uma agência bancária, já ordenado, foi:
$$
2, \ 3, \ 3, \ 5, \ 6, \ 7, \ 8, \ 9, \ 10, \ 12, \ 15, \ 30
$$

\begin{enumerate}[label=(\alph*)]
  \item Calcule $Q_1$, $Q_2$ (mediana) e $Q_3$.
  \item Calcule o $IQR$.
  \item Usando o critério $Q_3 + 1{,}5\times IQR$, o valor 30 minutos deve ser classificado como
    outlier? Mostre a conta.
\end{enumerate}
\end{questao}

\begin{questao}{4} \textbf{(Leitura de boxplot: comparação entre grupos)}

Uma consultoria compara o tempo (em dias) que propostas de crédito levam para ser aprovadas em
duas agências, A e B, por meio de boxplots lado a lado. A agência A tem caixa mais estreita e
mediana em 5 dias; a agência B tem caixa mais larga, mediana em 4 dias, mas com vários pontos
marcados como outliers acima de 20 dias.

\begin{enumerate}[label=(\alph*)]
  \item Qual agência tem processo de aprovação mais consistente (previsível)? Justifique usando o
    conceito de $IQR$.
  \item A agência B tem mediana menor que A. Isso significa que, em geral, B é mais rápida que A?
    O que os outliers acrescentam a essa conclusão?
  \item Que decisão gerencial diferente essa consultoria recomendaria para A e para B, com base
    apenas nesses dois boxplots?
\end{enumerate}
\end{questao}

\begin{questao}{5} \textbf{(Dedução: robustez da mediana)}

Considere um conjunto de dados $x_1 \le x_2 \le \cdots \le x_n$ ($n$ ímpar), com mediana
$x_{((n+1)/2)}$.

\begin{enumerate}[label=(\alph*)]
  \item Mostre que, se o maior valor $x_n$ for substituído por um valor arbitrariamente grande
    (mantendo-o ainda o maior da lista), a mediana \textbf{não se altera}. (Dica: pela definição, a
    mediana depende apenas da posição central da lista ordenada, não do valor exato dos elementos
    fora dessa posição.)
  \item Mostre, por meio de um contraexemplo numérico com $n=5$, que a \textbf{média} se altera nessa
    mesma situação.
  \item Com base em (a) e (b), explique por que a mediana é chamada de medida \textbf{robusta} a
    valores extremos, e a média não.
\end{enumerate}
\end{questao}

\bigskip
\noindent\textit{A questão a seguir não tem resposta única e não é cobrada em prova no mesmo
formato, é para discussão em grupo ou por escrito, antes da próxima aula.}

\begin{questao}{6 (discussão)} \textbf{(Sem resposta única)}

Procure (num relatório público, notícia ou site de estatísticas oficiais) um caso em que média e
mediana de uma mesma variável são divulgadas juntas. Se só uma das duas for divulgada, explique por
escrito qual você suspeita que esteja faltando e por quê a instituição pode ter escolhido divulgar
apenas uma delas.
\end{questao}


\bigskip
\section*{Exercícios complementares}

\begin{enumerate}[label=\textbf{C\arabic*.}, leftmargin=*]
\item O salário anual das secretárias de um escritório de advocacia é: R\$ 120.000, R\$ 60.000, R\$ 40.000, R\$ 40.000, R\$ 30.000, R\$ 30.000 e R\$ 30.000.
\begin{enumerate}[label=(\alph*)]
  \item Calcule a média, a moda e a mediana.
  \item Há diferença importante entre as três medidas? O que a explica?
  \item Qual medida representa melhor a tendência central desses salários? Justifique.
\end{enumerate}
\item Os preços semanais (em dólares) de uma marca de pizza congelada foram coletados em duas cidades dos EUA durante três anos:
\begin{center}
\begin{tabular}{l|ccccccc}
\toprule
Cidade & Média & Mediana & DP & Mín. & Máx. & $Q_1$ & $Q_3$ \\
\midrule
Dallas & 2,62 & 2,61 & 0,1559 & 2,21 & 3,05 & 2,51 & 2,72 \\
Chicago & 2,61 & 2,65 & 0,2122 & 1,55 & 3,03 & 2,53 & 2,74 \\
\bottomrule
\end{tabular}
\end{center}
\begin{enumerate}[label=(\alph*)]
  \item Calcule a distância interquartil dos preços em cada cidade.
  \item Há diferença na variabilidade dos preços entre as cidades? Compare o $IQR$ com a amplitude total e o desvio padrão.
  \item Em qual cidade o preço tende a ser mais barato? Justifique.
  \item Calcule os limites $Q_1 - 1{,}5\,IQR$ e $Q_3 + 1{,}5\,IQR$ para Chicago. O preço mínimo de 1,55 seria um valor atípico?
\end{enumerate}
\item Psicólogos acompanharam 20 voluntários com narcolepsia durante um dia e registraram quantas vezes cada um adormeceu:
\begin{center}
\begin{tabular}{c|ccccccc|c}
\toprule
Nº de cochilos & 1 & 2 & 3 & 4 & 5 & 6 & 7 & Total \\
\midrule
Freq. absoluta & 2 & 4 & 3 & 5 & 3 & 1 & 2 & 20 \\
\bottomrule
\end{tabular}
\end{center}
Encontre a moda, a média, a mediana e os quartis $Q_1$ e $Q_3$ (use o algoritmo de percentis da aula).
\item Nível educacional dos 39 funcionários administrativos de uma empresa: pós-graduação (3), superior completo (5), superior incompleto (11), médio completo (14), médio incompleto (6).
\begin{enumerate}[label=(\alph*)]
  \item Qual é o tipo da variável nível educacional?
  \item Forneça uma medida de tendência central adequada para essa variável. A média faria sentido?
\end{enumerate}
\item Considere a variável $X$ com a distribuição de frequências: $X = 9$ (1), $8$ (5), $0$ (1), $2$ (2), $6$ (1), $3$ (4), $4$ (1).
\begin{enumerate}[label=(\alph*)]
  \item Calcule a moda, a mediana e a média de $X$.
  \item Calcule $Q_1$ e $Q_3$ e verifique se há valores atípicos.
\end{enumerate}

\end{enumerate}

\bigskip
\needspace{12\baselineskip}
\section*{Aprofundamento: contas nacionais e dados reais}
\noindent\textit{Exercícios com dados oficiais, ligados à seção de contas nacionais do Capítulo 1 do livro. Use o algoritmo de percentis visto em aula.}

\begin{enumerate}[label=\textbf{A\arabic*.}, leftmargin=*]
\item \textbf{(Crescimento trimestral do PIB.)} Taxa de variação do PIB em relação ao trimestre imediatamente anterior, com ajuste sazonal (\%;
IBGE, Contas Nacionais Trimestrais), do 1º trimestre de 2023 ao 2º trimestre de 2026:
\begin{center}
1,3 \quad 0,8 \quad 0,1 \quad 0,4 \quad 0,7 \quad 1,7 \quad 0,9 \quad 0,1 \quad 1,3 \quad 0,4 \quad 0,1 \quad 0,2 \quad 1,1 \quad 0,5
\end{center}
\begin{enumerate}[label=(\alph*)]
  \item Calcule a média, a mediana, a moda, $Q_1$ e $Q_3$.
  \item Calcule a distância interquartílica e as cercas do boxplot. Há valores atípicos? Desenhe o boxplot.
  \item Uma notícia diz que ``a economia cresceu, em média, 0,7\% por trimestre''. Uma taxa trimestral de 0,7\% equivale a quanto ao ano, se mantida por quatro trimestres?
\end{enumerate}

\item \textbf{(PIB \emph{per capita} das Unidades da Federação.)} PIB \emph{per capita} das 27 UFs em 2021, em ordem crescente (R\$ mil; IBGE, Contas Regionais).
O menor é o do Maranhão e o maior, o do Distrito Federal:
\begin{center}
\small
17,5 \; 19,1 \; 19,5 \; 21,1 \; 22,2 \; 22,5 \; 22,7 \; 22,8 \; 22,9 \; 23,5 \; 23,6 \; 27,9 \; 30,0 \; 30,8 \\
32,0 \; 32,2 \; 37,4 \; 40,1 \; 45,4 \; 47,4 \; 50,1 \; 50,7 \; 54,4 \; 58,3 \; 58,4 \; 65,4 \; 92,7
\end{center}
\begin{enumerate}[label=(\alph*)]
  \item Calcule a média, a mediana, $Q_1$ e $Q_3$. Construa o boxplot e identifique valores atípicos.
  \item O PIB \emph{per capita} do Brasil em 2021 foi de R\$ 42,2 mil. Por que ele não coincide com a média calculada em (a)? Qual média
    ponderada reproduz o valor nacional?
  \item Qual medida de tendência central você usaria para descrever a UF ``típica''? Justifique.
\end{enumerate}

\end{enumerate}

\end{document}
